\section{Relation to raviolo vertex algebras}
\label{sec:raviolo}

\subsection{Restriction to time–slices}
Let $A=(A_{\mathsf{ch}},A_{\mathsf{top}})$ be an algebra over $C_\ast\!\bigl(W(\mathsf{SC}^{\mathrm{ch,top}})\bigr)$.
We define its restriction to a fixed time $t_0\in\mathbb R$ by evaluating on rectangles
$D\times I$ with $I$ a small interval centered at~$t_0$, and we identify objects along
the topological direction via the (contractible) $E_1$–homotopy.

\begin{definition}[Raviolo restriction]
\label{def:raviolo-restriction}
Fix $t_0\in\mathbb R$.  The \emph{raviolo factorization algebra} associated to~$A$ is the assignment
\[
  \mathsf V_{\mathrm{rav}}(D)
   \,:=\,
  \bigl( A_{\mathsf{ch}}(D\times I)\bigr)\big/\!\!\sim
  \qquad
  \text{for small disks $D\subset\mathbb C$ and intervals $I$ about $t_0$,}
\]
where $\sim$ identifies the two choices of approaching $t_0$ ("from above" and "from below")
using the $E_1$–coherence data.
\end{definition}

\begin{proposition}[\ClaimStatusProvedHere]\label{prop:raviolo-VA}
The restriction $\mathsf V_{\mathrm{rav}}$ carries the structure of a (dg) raviolo vertex algebra in the sense of \cite{GW23,AKY24}.
\end{proposition}
\begin{proof}
We verify the raviolo vertex algebra axioms by constructing the structure maps from the $\OHT$-algebra data and checking each axiom.

\medskip\noindent\textbf{Step 1: Two-sheeted structure from time-slice restriction.}
By Proposition~\ref{prop:SC-raviolo}, the time-slice restriction of the $\OHT$-algebra $(A_{\mathsf{ch}}, A_{\mathsf{top}})$ to a fixed $t_0\in\R$ produces a raviolo factorization algebra $\mathsf{V}_{\mathrm{rav}}$ on $\C$.  The raviolo geometry arises as follows: the two sheets $D_+$ and $D_-$ correspond to the two half-intervals $I_+ = (t_0, t_0+\epsilon)$ and $I_- = (t_0-\epsilon, t_0)$, representing the ``future'' and ``past'' of the time-slice.  The punctured disk $D^\times$ corresponds to the full interval $I = I_-\cup I_+$ minus the point $\{t_0\}$, and the $E_1$-homotopy data provides the gluing $D_+\cup_{D^\times} D_-$.  The state space is $V := A_{\mathsf{ch}}(D\times I)$ for a small disk $D\subset\C$.

\medskip\noindent\textbf{Step 2: Vertex algebra operations.}
The vertex algebra structure maps on $\mathsf{V}_{\mathrm{rav}}$ are obtained by restricting the $A_\infty$ chiral operations $m_k$ to the time-slice.  The state-field correspondence $Y(a,z): V\to V((z))$ is given by the binary operation:
\[
  Y(a,z)\,b \;=\; m_2(a,b)\big|_{\lambda = z^{-1}},
\]
where the spectral parameter $\lambda$ is related to the holomorphic coordinate difference by the standard identification $\lambda \leftrightarrow z^{-1}$ (Fourier--Laplace duality between the spectral and position descriptions).  The vacuum $|0\rangle$ is the unit of $A_{\mathsf{ch}}$ (the identity observable), and the translation operator $T = \partial$ acts by the holomorphic derivative.

\medskip\noindent\textbf{Step 3: Axiom verification.}
\begin{enumerate}[label=(\roman*)]
\item \emph{Vacuum axiom}: $Y(|0\rangle, z) = \mathrm{id}_V$ follows from the $A_\infty$ unit axiom $m_2(\mathbf{1},a) = a$.  The creation property $Y(a,z)|\mathbf{0}\rangle\big|_{z=0} = a$ follows from the unit axiom $m_2(a,\mathbf{1}) = a$.

\item \emph{Translation covariance}: $[T, Y(a,z)] = \partial_z Y(a,z)$ is precisely the sesquilinearity identity $m_2(\partial a, b) = -\lambda\, m_2(a,b)$ rewritten in position space, where $\partial_z$ corresponds to multiplication by $-\lambda$ under Fourier--Laplace.

\item \emph{Locality}: For $a,b\in V$, the commutator $[Y(a,z), Y(b,w)]$ is a distribution supported on $z=w$, annihilated by $(z-w)^N$ for sufficiently large $N$.  Since $m_2(a,b)\in A((\lambda))$ has only finitely many singular terms in $\lambda^{-1}$, the OPE has finite polar order.

\item \emph{Raviolo structure}: The two sheets $D_\pm$ carry the vertex algebra operations from (ii)--(iii), and the $E_1$-coherence data provides the homotopy connecting the two sheets.  Concretely, the $A_\infty$ operations $m_k$ for $k\ge 3$ encode the higher homotopies of the raviolo: they measure the failure of the binary vertex algebra structure to be strictly associative, with the homotopy-coherent corrections organized by the topological $E_1$-direction.  The passage from $\OHT$ to a raviolo VA is the statement that the closed-color operations (holomorphic, governed by $\FM_k(\C)$) restrict to vertex algebra operations on each sheet, while the open-color operations (topological, governed by $E_1(m)\cong\Conf_m(\R)$) provide the gluing data between sheets.
\end{enumerate}
This establishes the (dg) raviolo vertex algebra structure on $\mathsf{V}_{\mathrm{rav}}$ in the sense of \cite{GW23, AKY24}.
\end{proof}

\begin{proposition}[Swiss–cheese $\Rightarrow$ raviolo; \ClaimStatusProvedHere]
\label{prop:SC-raviolo}
Restriction of a $C_\ast\!\bigl(W(\mathsf{SC}^{\mathrm{ch,top}})\bigr)$–algebra to a time–slice yields
a (dg) raviolo factorization algebra on~$\mathbb C$.  Its product and OPE maps are
induced by those of $A$ via the canonical $E_1$–homotopies.
\end{proposition}

\begin{proof}
By the local constancy in the topological direction, an interval $I$ about~$t_0$ can be
retracted to a point using $E_1(1)$, which is contractible.  The contractibility of $E_1(1)$ means that the topological operations on a single interval are homotopically trivial. Restricting to a time-slice collapses the interval to a point, but the two endpoints (incoming and outgoing) retain distinct identities, producing the two sheets $D_+$ and $D_-$ of the raviolo. The open disk $D^\times$ parametrizes the interior of the interval, and the $E_1$-homotopy provides the gluing data $D_+ \cup_{D^\times} D_-$.  Compatibility with disjoint unions
follows from the swiss–cheese composition rules on rectangles.
\end{proof}

\begin{remark}[Čech/Thom–Sullivan model]
\label{rmk:TS-model}
Equivalently, $\mathsf V_{\mathrm{rav}}$ can be computed by a Čech resolution of
the cover $\{D_+,D_-,D^\times\to D_+\cup_{D^\times}D_-\}$ and then passing to
Thom–Sullivan chains; this is the derived section model used in
\cite{GW23,AKY24}.
\end{remark}

\begin{remark}[HCA structure on the raviolo \v{C}ech complex]
\label{rmk:hca-raviolo-cech}
\index{homotopy chiral algebra!raviolo}
The three-element cover $\mathcal{U} = \{D_+, D_-, D^\times\}$
of the raviolo produces a cosimplicial Lie algebra in the
pseudo-tensor category: sections of the boundary chiral algebra
$\cA_\partial$ on $D_+$, $D_-$, $D^\times$, and their
pairwise/triple intersections, with coface maps given by
restriction.  By the Malikov--Schechtman theorem~\cite{MS24}
(Theorem~3.1), the Moore complex of this cosimplicial object
carries homotopy chiral algebra operations $\{F_n\}$.

Concretely, the normalized Moore complex is the three-term complex
\[
(NM\check{C})^0 = \cA_\partial(D_+),
\quad
(NM\check{C})^1
  = \ker\bigl(\cA_\partial(D_+ \cap D^\times) \to \cA_\partial(D^\times)\bigr),
\quad
(NM\check{C})^2 = 0.
\]
The vanishing in degree~$2$ follows from the Moore normalization: the
triple intersection $D_+ \cap D_- \cap D^\times = D^\times$ is nonempty,
but the normalized Moore condition (intersection of kernels of all
degeneracy maps) kills the sections supported on~$D^\times$, reducing
the complex to two terms.  Since $\check{C}^2 = 0$, the \v{C}ech MC element
(from the strict \v{C}ech HCA package of Volume~I) is \emph{strict} for
every boundary chiral algebra:
\[
\Phi_{\cA_\partial}
= (d_M,\; F_2,\; 0,\; 0,\; \ldots)
\in \operatorname{MC}\!\bigl(\mathfrak{g}^{\check{C}}_{\cA_\partial}\bigr),
\]
where $F_2$ is the chiral bracket on sections restricted to
$D_+ \cap D^\times$.  Explicitly, $F_2$ acts on
$(M\check{C})^0 \times (M\check{C})^1 \to (M\check{C})^1$ by
restricting a section $a \in \cA_\partial(D_+)$ to
$D_+ \cap D^\times$ and applying the chiral bracket:
\[
F_2(a, \bar{b})
  \;=\; \overline{\mu^{\mathrm{ch}}(a|_{D_+ \cap D^\times},\, b)},
\qquad
a \in (M\check{C})^0,\;\;
\bar{b} \in (M\check{C})^1,
\]
where $\bar{b}$ denotes the class of $b$ in
$\cA_\partial(D_+ \cap D^\times)/\cA_\partial(D^\times)$ and
$\mu^{\mathrm{ch}}$ is the chiral bracket of $\cA_\partial$.
The vanishing of $F_3$ follows from
the two-element strictness proposition of Volume~I: the
Jacobiator targets $\check{C}^2 = 0$.

The Thom--Sullivan passage from \v{C}ech to de~Rham preserves
this strictness, and the raviolo PVA of
Theorem~\ref{thm:raviolo-PVA} is computed by a strict model.
The non-trivial $A_\infty$ operations $m_3$ on the bar side
(which can be nonzero for interacting theories) enter only
through the \emph{comparison morphism} $\Phi^*$, not through
the \v{C}ech model itself.  This is the descent-theoretic
explanation for why the PVA is strict despite the bar
$A_\infty$ structure being non-formal: the raviolo cover has
no triple intersections, so no room for Jacobiators.
\end{remark}

\subsection{Cohomology and shifted Poisson vertex structure}
\label{subsec:raviolo-PVA}
The cohomological vertex--Poisson algebra structure of
Section~\ref{sec:Ainfty-to-PVA} applies: the symmetric regular part of~$m_2$ gives the
(commutative, associative) product on cohomology and the antisymmetric singular part gives a
degree~$-1$ $\lambda$--bracket.

\begin{theorem}[PVA structure on raviolo cohomology; \ClaimStatusProvedHere]
\label{thm:raviolo-PVA}
Let $A$ be a $C_\ast\!\bigl(W(\mathsf{SC}^{\mathrm{ch,top}})\bigr)$–algebra
and let $\mathsf V_{\mathrm{rav}}$ be its raviolo restriction.
Then $H^\bullet(\mathsf V_{\mathrm{rav}},Q)$ carries a $(-1)$–shifted
Poisson vertex algebra structure $(\cdot,\{\ \_ \})$:
\begin{enumerate}[label=(\alph*)]
\item $(\cdot)$ is a commutative, associative product induced by the symmetric regular part of $m_2$;
\item $\{\,\cdot{}_\lambda \cdot\,\}$ is a degree~$-1$ bracket induced by the antisymmetric singular part of $m_2$;
\item sesquilinearity, skewsymmetry, the Jacobi identity, and Leibniz rule hold on cohomology.
\end{enumerate}
\end{theorem}

\begin{proof}
The PVA structure on $H^\bullet(\mathsf{V}_{\mathrm{rav}}, Q)$ is extracted from the chain-level $A_\infty$ chiral algebra via the regular/singular decomposition of $m_2$.  We verify each axiom in turn.

\medskip\noindent\textbf{Step 1: Commutative product from the regular part.}
Write $m_2(a,b) = m_2^{\mathrm{reg}}(a,b) + m_2^{\mathrm{sing}}(a,b)$ where $m_2^{\mathrm{reg}}(a,b)\in A[[\lambda]]$ and $m_2^{\mathrm{sing}}(a,b)\in \lambda^{-1}A[[\lambda^{-1}]]$.  The product on cohomology is defined by the symmetrized regular part:
\[
  [a]\cdot [b] \;:=\; \tfrac{1}{2}\Bigl(\bigl[m_2^{\mathrm{reg}}(a,b)\big|_{\lambda=0}\bigr] + (-1)^{|a||b|}\bigl[m_2^{\mathrm{reg}}(b,a)\big|_{\lambda=0}\bigr]\Bigr).
\]
The time-slice restriction maps the $\OHT$-algebra product to the raviolo product: on the time-slice, the interval $I$ collapses and the $\lambda=0$ evaluation corresponds to the coincident-time limit $t_1=t_2$, which is well-defined because the regular part has no singularity.  Commutativity is manifest from the symmetrization.  Associativity follows from the $n=3$ $A_\infty$ identity (Theorem~\ref{thm:assoc}): the obstruction to associativity is $Q(m_3^{\mathrm{reg}})$, which vanishes on $Q$-closed inputs.

\medskip\noindent\textbf{Step 2: $\lambda$-bracket from the singular part.}
The $\lambda$-bracket on cohomology is defined by the antisymmetrized singular part:
\[
  \{[a]{}_\lambda [b]\} \;:=\; \bigl[m_2^{\mathrm{sing}}(a,b)\bigr] - (-1)^{(|a|+1)(|b|+1)}\bigl[m_2^{\mathrm{sing}}(b,a)\big|_{\lambda\mapsto -\lambda-\partial}\bigr].
\]
On the raviolo, the singular part $m_2^{\mathrm{sing}}$ encodes the OPE poles: the coefficient of $\lambda^{-n-1}$ in $m_2^{\mathrm{sing}}(a,b)$ corresponds to the $(n)$-th product $a_{(n)}b$ in vertex algebra language.  The time-slice restriction preserves these poles because singularities in the holomorphic direction survive the collapse of the topological interval.

\medskip\noindent\textbf{Step 3: Sesquilinearity.}
The chain-level sesquilinearity axioms $m_2(\partial a, b) = -\lambda\, m_2(a,b)$ and $m_2(a, \partial b) = (\lambda+\partial)\,m_2(a,b)$ (Equations~\ref{eq:sesqui-left}--\ref{eq:sesqui-right}) descend directly to cohomology.  Since $Q=m_1$ commutes with $\partial$ (the translation operator acts on the holomorphic coordinate, not the cohomological degree), and since $Q$ does not involve the spectral parameter $\lambda$, the sesquilinearity identities pass through the Laurent decomposition and the projection to $Q$-cohomology without modification:
\[
  \{\partial[a]{}_\lambda [b]\} = -\lambda\,\{[a]{}_\lambda [b]\}, \qquad
  \{[a]{}_\lambda \partial[b]\} = (\lambda+\partial)\,\{[a]{}_\lambda [b]\}.
\]

\medskip\noindent\textbf{Step 4: Skew-symmetry.}
The antisymmetrization in Step~2 gives, by the exchange $\lambda\mapsto -\lambda-\partial$ and the Koszul sign $(-1)^{(|a|+1)(|b|+1)}$:
\[
  \{[a]{}_\lambda [b]\} \;=\; -(-1)^{(|a|+1)(|b|+1)}\,\{[b]{}_{-\lambda-\partial} [a]\}.
\]
This is the standard skew-symmetry for a $(-1)$-shifted $\lambda$-bracket, matching the degree convention $|m_2|=1-2=-1$.

\medskip\noindent\textbf{Step 5: Jacobi identity.}
The $n=3$ $A_\infty$ identity, projected to the singular part in all three pairwise channels, yields the $\lambda$-Jacobi identity (Theorem~\ref{thm:Jacobi-PVA}).  On the raviolo, this has the following geometric interpretation: the three boundary strata $D_{12}$, $D_{13}$, $D_{23}$ of $\FM_3(\C)$ contribute the three terms of the Jacobi identity, corresponding to the three pairwise collisions of holomorphic insertion points on the time-slice.  The total collision stratum $D_{123}$ contributes a $Q$-exact term $Q(m_3^{\mathrm{sing}})$ that vanishes on cohomology classes.  The Arnold relation $d\log(z_1-z_2)\wedge d\log(z_2-z_3) + \text{cyclic} = 0$ ensures cancellation at codimension-2 corners, so the Stokes argument closes.

\medskip\noindent\textbf{Step 6: Leibniz rule.}
The $n=3$ $A_\infty$ identity with mixed projection --- outer $m_2$ regular (product), inner $m_2$ singular (bracket) --- yields the Leibniz rule (Theorem~\ref{thm:Leibniz-PVA}):
\[
  \{[a]{}_\lambda ([b]\cdot [c])\}
  \;=\; \{[a]{}_\lambda [b]\}\cdot [c]
  \;+\; (-1)^{(|a|+1)|b|}\,[b]\cdot\{[a]{}_\lambda [c]\}.
\]
This mixed projection is well-defined because the regular/singular decomposition is multiplicative: the spectral substitution principle (Remark~\ref{rmk:spectral-substitution}) governs how inner spectral parameters compose with outer ones, and the regular part of a composition with an inner singularity isolates the bracket-then-multiply structure.

\medskip\noindent\textbf{Step 7: Degree shift.}
The binary operation $m_2$ has intrinsic degree $|m_2|=1-2=-1$.  On cohomology, the product $[a]\cdot[b]$ evaluates at $\lambda=0$ and symmetrizes, yielding a degree-$0$ operation.  The bracket $\{[a]{}_\lambda [b]\}$ retains the spectral parameter, and extracting the coefficient of $\lambda^n$ preserves the degree shift of $-1$ from $m_2$.  As a binary operation on the underlying graded vector space, the bracket is therefore a map $V\otimes V\to V$ of degree $+1$ (absorbing the $-1$ from $|m_2|$ and the $+2$ from the two desuspensions implicit in the $\lambda$-bracket formalism), consistent with a $(-1)$-shifted Poisson vertex algebra.

\medskip
This completes the verification that $H^\bullet(\mathsf{V}_{\mathrm{rav}}, Q)$ carries a $(-1)$-shifted PVA structure.
\end{proof}

\begin{corollary}[Compatibility with Oh–Yagi and CDG; \ClaimStatusConditional]
\label{cor:compatibility-OY-CDG}
For HT twists of $3$d~$\mathcal N=2$ theories satisfying Theorem~\ref{thm:physics-bridge},
the $(-1)$–shifted Poisson vertex structure on $H^\bullet(\mathsf V_{\mathrm{rav}},Q)$
agrees with the secondary bracket constructed in
\cite{OY20} and \cite{CDG20}.
\end{corollary}

\begin{proof}
Both brackets are computed by the same residue of the $m_2$ singular kernel,
and both products arise from the regular part of the binary operation.
\end{proof}

\subsection{Higgs/Coulomb branch checks}
\label{subsec:Higgs-Coulomb}
For $3$d~$\mathcal N=4$ theories, the raviolo vertex algebra specializes to the
superconformal raviolo vertex algebras of
\cite{GRW21}, where $H^\bullet(\mathsf V_{\mathrm{rav}},Q)$
recovers ring of functions on Higgs and Coulomb branches.
The $2$–shifted Poisson structures observed therein descend to the $(-1)$–shifted
vertex Poisson structure described in Theorem~\ref{thm:raviolo-PVA}.
\begin{proposition}[Higgs/Coulomb compatibility; \ClaimStatusConditional]
\label{prop:Higgs-Coulomb-compat}
Let $\mathcal T$ be a $3$d~$\mathcal N=4$ theory with an HT twist satisfying Theorem~\ref{thm:physics-bridge}.
Then $H^\bullet(\mathsf V_{\mathrm{rav}},Q)$
is simultaneously a quantization of the Higgs and Coulomb branch Poisson algebras, compatible with the
raviolo PVA structure.  In particular, restriction to holomorphic currents produces the
expected $\widehat{\mathfrak g}$–type current algebras on the boundary.
\end{proposition}
\begin{proof}[Sketch]
This is verified in examples by direct calculation:
free multiplets, abelian gauge theories with Chern–Simons coupling, and theories with Landau–Ginzburg
superpotentials.  The compatibility of OPEs with the raviolo factorization encodes the algebraic structure and
identifies the shifted Poisson bracket with the secondary operation of \cite{OY20}.
\end{proof}
\subsection{Examples and consistency checks}
\label{subsec:raviolo-examples-2}
We briefly recall the behavior in the three model families analysed in Section~\ref{sec:examples}.
\paragraph{\textbf{Free chiral multiplet.}}
The singular part of $m_2$ vanishes in cohomology; the regular part produces a purely commutative vertex algebra.
This matches the standard free–field realization on the raviolo.
\paragraph{\textbf{Landau–Ginzburg model.}}
With a cubic superpotential, the first nontrivial $m_3$ may appear, but its contribution vanishes on cohomology
by the usual $A_\infty$ argument; the singular residue produces a nontrivial degree $-1$ bracket on $H^\bullet$,
consistent with the calculations in \cite{GRW21} and Section~\ref{sec:examples}.

\paragraph{\textbf{Abelian Chern–Simons boundary.}}
On the boundary, the raviolo vertex algebra contains $\widehat{\mathfrak u}(1)$ currents; the singular part of $m_2$
computes the level and the bracket on cohomology agrees with the current algebra OPE.
The bulk–to–boundary map (Section~\ref{sec:bulk-boundary}) intertwines the bulk PVA with the boundary VOA.

\subsection{The essential principle}
\label{sec:bridge-principle}

The bar complex $\barB^{\mathrm{ch}}(\cA) = T^c(s^{-1}\bar\cA)$ carries a differential $d_{\barB}$ from OPE residues on $\FM_k(\C)$ and a coproduct~$\Delta$ from deconcatenation on $\Conf_k(\R)$.  A bar element of degree~$k$ is parametrised by $\FM_k(\C) \times \Conf_k(\R)$ --- the operation spaces of~$\SCchtop$ (Theorem~\ref{thm:steinberg-presentation}).

\begin{definition}[Logarithmic $\SCchtop$-algebra]
\label{def:log-SC-algebra}
A \emph{logarithmic $\SCchtop$-algebra} is a $C_\ast(W(\SCchtop))$-algebra $(\cA_{\mathrm{ch}}, \cA_{\mathrm{top}})$ whose closed-colour $A_\infty$ operations $\{m_k\}_{k \ge 1}$ are defined by weight forms
\[
  \omega_k \;=\; \omega_k^{\mathrm{hol}} \otimes \omega_k^{\mathrm{top}},
  \qquad
  \omega_k^{\mathrm{hol}} \in \Omega^\bullet_{\log}(\FM_k(\C)),
  \quad
  \omega_k^{\mathrm{top}} \in C_\bullet(\Conf_k(\R)),
\]
satisfying the product decomposition: weight forms factor as holomorphic logarithmic forms on the Fulton--MacPherson compactification tensored with topological chains on ordered configurations.
\end{definition}

\begin{theorem}[FM calculus for logarithmic $\SCchtop$-algebras; \ClaimStatusProvedHere]
\label{thm:FM-calculus}
Let $(\cA_{\mathrm{ch}}, \cA_{\mathrm{top}})$ be a logarithmic $\SCchtop$-algebra.  Then:
\begin{enumerate}[label=(\alph*)]
\item \textbf{Logarithmic forms}: The weight forms $\omega_k^{\mathrm{hol}} \in \Omega^\bullet_{\log}(\FM_k(\C))$ have prescribed pole orders at each boundary divisor, determined by the $A_\infty$ structure maps.

\item \textbf{Factorisation at boundaries}: On each divisor $D_S \subset \partial \FM_k(\C)$, the restriction $\omega_k|_{D_S}$ factors compatibly with the FM boundary structure:
\[
\FM_k(\C)|_{D_S} \cong \FM_{|S|}(\C) \times \prod_{i \in S} \FM_{|C_i|}(\C),
\]
where $C_i$ are the connected components of the cluster $S$.

\item \textbf{Arnold relations}: Logarithmic forms satisfy the classical Arnold--Orlik--Solomon relations at codimension-$2$ corners $D_S \cap D_T$, ensuring compatibility of iterated residues.

\item \textbf{Stokes exactness}: The exterior derivative $d$ satisfies Stokes' theorem on manifolds with corners, and corner contributions cancel via Arnold relations (Appendix~\ref{app:FM_Stokes}).
\end{enumerate}
Consequently, the $A_\infty$ relations
$\sum_{i+j=n+1} m_i \circ (\mathrm{id}^{\otimes p} \otimes m_j \otimes \mathrm{id}^{\otimes q}) = 0$
hold by Stokes' theorem on $\FM_k(\C)$, with codimension-$2$ corner contributions cancelling via Arnold--Orlik--Solomon relations.
\end{theorem}

\begin{proof}
Each assertion is a theorem of configuration space geometry, independent of any physical realisation.  Parts~(a)--(c) follow from the theory of logarithmic forms on wonderful compactifications (cf.~\cite{FM94,DeConciniProcesi}); part~(d) is proved in Appendix~\ref{app:FM_Stokes}.
\end{proof}

\begin{theorem}[Bridge from physics to algebra; \ClaimStatusProvedHere]
\label{thm:physics-bridge}
Let $\mathcal{T}$ be a 3d holomorphic--topological quantum field theory on $\R_t \times \C_z$ with:
\begin{enumerate}[label=(\alph*)]
\item \textbf{BV data and HT gauge fixing}: A $(-1)$-symplectic graded field space $\mathcal{F}$ with local action $S \in \mathcal{O}(\mathcal{F})$ satisfying $\{S,S\}_{\mathrm{BV}} = 0$, gauge-fixed in HT gauge with $Q_{\mathrm{BRST}}^2 = 0$.

\item \textbf{One-loop finiteness}: In the holomorphic--topological gauge, one-loop Feynman integrals converge without UV divergences (the Gwilliam--Rabinovich--Williams mechanism~\cite{GRW21} for $d' = 1$).

\item \textbf{Polynomial interactions}: The interaction vertices are polynomial in the fields.
\end{enumerate}
Then the BV--BRST complex $\Obs_{\mathrm{bulk}}$ of $\mathcal{T}$, equipped with configuration-space $A_\infty$ operations, is a logarithmic $\SCchtop$-algebra in the sense of Definition~\ref{def:log-SC-algebra}.
\end{theorem}

\begin{proof}
\textbf{Meromorphicity.}  The gauge-fixed propagator is the Green's function of $Q = \dbar + d_t$, computed as the tensor product of the Cauchy kernel $1/(2\pi z)$ (solving $\dbar G_\C = \delta$) with the Heaviside distribution~$\Theta(t)$ or Dirac~$\delta(t)$ (solving $d_t G_\R = \delta$).  This is meromorphic in $\C$ with a simple pole at $z = 0$ and distributionally supported or exponentially decaying in~$\R$.

\textbf{Product factorisation.}
We prove that the Feynman weight forms factor as
$\omega_\Gamma = \omega_\Gamma^{\mathrm{hol}} \otimes \omega_\Gamma^{\mathrm{top}}$
with $\omega_\Gamma^{\mathrm{hol}} \in \Omega^\bullet_{\log}(\FM_k(\C))$ and
$\omega_\Gamma^{\mathrm{top}} \in C_\bullet(\Conf_k(\R))$.
The argument has four steps: propagator factorisation, vertex factorisation,
integrand factorisation, and integration factorisation.

\emph{Step (i): Propagator factorisation.}
The BRST operator splits as $Q = \dbar_\C + d_t$ with
$\dbar_\C$ acting on $\C_z$ and $d_t$ acting on $\R_t$.
Since $Q$ is a sum of operators on different factors,
the gauge-fixed Green's function $G$ solving $QG = \delta^{(3)}$
decomposes as a tensor product of the fundamental solutions
on each factor:
\[
  G(z_{ij};\, t_{ij})
  \;=\;
  K_\C(z_{ij}) \,\otimes\, H_\R(t_{ij}),
  \qquad
  K_\C(z) = \frac{1}{2\pi z},
  \quad
  H_\R(t) = \Theta(t)\text{ or }\delta(t),
\]
where $z_{ij} = z_i - z_j$, $t_{ij} = t_i - t_j$,
$\dbar K_\C = \delta_\C$, and $d_t H_\R = \delta_\R$.
The Cauchy kernel $K_\C$ is meromorphic in $z$ with a simple
pole at the origin; $H_\R$ is distributional on~$\R$.
Each propagator therefore splits into a holomorphic factor
(depending only on $\C$-coordinates) and a topological factor
(depending only on $\R$-coordinates).

\emph{Step (ii): Vertex factorisation.}
The interaction vertices of the theory are polynomial in the
fields and their derivatives (item~(c)).  In the HT
gauge, the fields decompose as sections of bundles on
$\C \times \R$.  Since the interaction Lagrangian is a
polynomial in $\Phi$ and its derivatives
$\partial_z^a \partial_t^b \Phi$, each monomial in the
Taylor expansion of the vertex factor at a point
$(z_v, t_v) \in \C \times \R$ is a product of a holomorphic
function of the $z$-derivatives and a function of the
$t$-derivatives:
\[
  V_v
  \;=\;
  \sum_{\alpha} V_v^{\mathrm{hol},\alpha}(z_v,\partial_z)
  \;\otimes\;
  V_v^{\mathrm{top},\alpha}(t_v,\partial_t).
\]
This finite sum structure arises because the interaction is
polynomial: each monomial in the fields has a definite number
of $z$- and $t$-derivatives, and no mixed-direction couplings
beyond those dictated by the product structure $\C \times \R$.

\emph{Step (iii): Integrand factorisation.}
The weight form for a Feynman graph~$\Gamma$ with edges~$E(\Gamma)$
and internal vertices~$V_{\mathrm{int}}(\Gamma)$ is
(Definition~\ref{def:weight_forms})
\[
  \omega_\Gamma
  \;=\;
  \Bigl(\prod_{e \in E(\Gamma)} G_e\Bigr)
  \cdot
  \Bigl(\prod_{v \in V_{\mathrm{int}}(\Gamma)} V_v\Bigr)
  \cdot d\mu.
\]
By Steps~(i) and~(ii), each propagator factor $G_e$
and each vertex factor $V_v$ is a (finite sum of) tensor
products of a $\C$-factor and an $\R$-factor.
The measure decomposes as
$d\mu = d\mu_\C \otimes d\mu_\R$
with $d\mu_\C = \prod d^2 z_v$ and
$d\mu_\R = \prod dt_v$.
Since the tensor product of finite sums of pure tensors
is again a finite sum of pure tensors, the full integrand
factors:
\[
  \omega_\Gamma
  \;=\;
  \sum_\alpha
  \omega_{\Gamma,\alpha}^{\mathrm{hol}}(z_1,\ldots,z_k;\,z_{v_1},\ldots)
  \;\otimes\;
  \omega_{\Gamma,\alpha}^{\mathrm{top}}(t_1,\ldots,t_k;\,t_{v_1},\ldots),
\]
where each $\omega_{\Gamma,\alpha}^{\mathrm{hol}}$ depends only on the
holomorphic coordinates and each
$\omega_{\Gamma,\alpha}^{\mathrm{top}}$ depends only on the
topological coordinates.

\emph{Step (iv): Integration factorisation.}
The $A_\infty$ operation $m_k$ is obtained by integrating
$\omega_\Gamma$ over the positions of the internal vertices
and summing over graphs (Convention~\ref{const:regularized_mk}).
The integration domain for internal vertices decomposes as
$\C^{|V_{\mathrm{int}}|} \times \R^{|V_{\mathrm{int}}|}$.
By Fubini's theorem (applicable because one-loop finiteness
(item~(b)) guarantees absolute convergence of the
integrals after FM compactification), integration over
internal vertex positions respects the tensor decomposition:
\[
  \int_{\C^{|V_{\mathrm{int}}|} \times \R^{|V_{\mathrm{int}}|}}
  \omega_\Gamma
  \;=\;
  \sum_\alpha
  \biggl(\int_{\C^{|V_{\mathrm{int}}|}}
  \omega_{\Gamma,\alpha}^{\mathrm{hol}}\biggr)
  \;\otimes\;
  \biggl(\int_{\R^{|V_{\mathrm{int}}|}}
  \omega_{\Gamma,\alpha}^{\mathrm{top}}\biggr).
\]
The holomorphic integrals produce meromorphic functions of the
external $z$-coordinates with poles at collision loci;
on~$\FM_k(\C)$ these extend to logarithmic forms
$\omega_k^{\mathrm{hol}} \in \Omega^\bullet_{\log}(\FM_k(\C))$
(by the logarithmic extension of $dz/(z_i - z_j) = d\log(z_i - z_j)$
at each boundary divisor).
The topological integrals, combined with the Heaviside
time-ordering from $H_\R$, produce chains
$\omega_k^{\mathrm{top}} \in C_\bullet(\Conf_k(\R))$ supported
on the ordered simplices $\{t_1 > \cdots > t_k\}$.
The resulting weight form is
\[
  \omega_k
  \;=\;
  \omega_k^{\mathrm{hol}} \otimes \omega_k^{\mathrm{top}}
  \;\in\;
  \Omega^\bullet_{\log}(\FM_k(\C)) \otimes C_\bullet(\Conf_k(\R)),
\]
as required by Definition~\ref{def:log-SC-algebra}.

\smallskip
\noindent\emph{Where the bridge-theorem inputs enter.}
The splitting $Q = \dbar + d_t$ enters in Step~(i)
(propagator factorisation).  Polynomial interactions
enter in Step~(ii) (ensuring the vertex factors are
finite sums of pure tensors; a non-polynomial interaction
such as $e^{\Phi(z,t)}$ could produce an infinite series
whose $\C$- and $\R$-dependence is entangled via
convergence issues).  One-loop finiteness enters in
Step~(iv) (justifying the application of Fubini's theorem
to the FM-compactified integrals).

\textbf{$\SCchtop$ structure.}  By the recognition theorem (Theorem~\ref{thm:recognition-SC}), the resulting holomorphic--topological prefactorisation algebra on $\C \times \R$ is a $C_\ast(W(\SCchtop))$-algebra.
\end{proof}

\begin{example}[Free Chiral Multiplet Propagator]
\label{ex:free_propagator}
For the free chiral multiplet $\Phi$ with action
\[
S_{\mathrm{free}} = \int_{\R \times \C} \bar{\Phi} \wedge (\partial_t + \bar{\partial}) \Phi,
\]
the propagator $K = \delta(t_1 - t_2) / (2\pi(z_1 - z_2))$ solves $(\partial_t + \bar{\partial}) K = \delta^{(3)}$.  This exhibits the product decomposition: meromorphic in $z$ with a simple pole, distributional $\delta(t)$ in the topological direction, and translation-invariant.
\end{example}

\begin{remark}[Physical realisations]
\label{rem:physical-realisations}
Theorem~\ref{thm:physics-bridge} is verified by holomorphic twists of 3d $\mathcal{N}=2$ theories~\cite{CDG20} and by HT theories with $d'=1$~\cite{GRW21, GKW25}.  The worked examples of Part~\textup{IV} are explicit verifications.
\end{remark}

\subsubsection*{II. Chain-level structure (Section \ref{sec:BV_construction})}

\begin{theorem}[Chain-level $A_\infty$ chiral structure; \ClaimStatusProvedHere]
\label{thm:chain_level_structure}
Let $(\cA_{\mathrm{ch}}, \cA_{\mathrm{top}})$ be a logarithmic $\SCchtop$-algebra.  Then $\cA_{\mathrm{ch}}$, equipped with operations $m_k$ from integration of weight forms over $\Conf_k(\R) \times \FM_k(\C)$, is an $A_\infty$ chiral algebra.  The $A_\infty$ relations follow from Stokes' theorem on $\FM_k(\C)$ with corner cancellation via Arnold--Orlik--Solomon (Theorem~\ref{thm:stokes_arnold}).
\end{theorem}

\subsubsection*{III. PVA descent (Section \ref{sec:PVA_descent})}

\begin{theorem}[PVA descent; \ClaimStatusProvedHere]
\label{thm:PVA-descent-roadmap}
For any logarithmic $\SCchtop$-algebra, $H^\bullet(\cA_{\mathrm{ch}},Q)$ is a $(-1)$-shifted Poisson vertex algebra: the $\lambda$-bracket descends from $m_2$, and $m_{k \geq 3}$ vanish on cohomology via homotopies $h_{k-1}$ from topological contraction of $\Conf_k^{<}(\R)$ (Theorem~\ref{thm:cohomology_PVA}).
\end{theorem}

\subsubsection*{IV. Chiral Hochschild and bar--cobar duality (Sections \ref{sec:chiral_hochschild_expanded}--\ref{sec:bar_cobar})}

\begin{theorem}[Bulk as chiral Hochschild; \ClaimStatusProvedHere]
\label{thm:bulk-hochschild-roadmap}
$C^\bullet_{\mathrm{ch}}(A_\partial) \simeq \A_{\text{bulk}}$ as $A_\infty$ chiral algebras (Theorem~\ref{thm:bulk_hochschild}).  The filtered bar--cobar adjunction extending~\cite{GLZ22} has differential $d_{\overline{B}}$ from FM residue sums with $d_{\overline{B}}^2 = 0$ by Arnold relations (Theorem~\ref{thm:bar_cobar_adjunction}).
\end{theorem}

\phantomsection\label{sec:bridge-inputs-summary}

\subsection{The $d'=1$ formality gap}

Let $d'$ denote the number of topological directions. The Gaiotto--Kulp--Wu formality theorem \cite{GKW25} establishes:

\begin{itemize}
\item \textbf{$d' \geq 2$}: Higher $A_\infty$ operations beyond $m_2$ stabilize---quantum corrections vanish after one loop and the theory is formal.

\item \textbf{$d' = 1$}: Formality fails. Higher $A_\infty$ operations $m_k$ can persist at all loop orders.
\end{itemize}

Three-dimensional HT theories on $\R \times \C$ have $d'=1$. When and why higher operations vanish in this setting remains open.

\subsection{Organization}

Section~\ref{sec:HT-operad} constructs the HT operad $\OHT$ and defines $A_\infty$ chiral algebras.
Section~\ref{sec:BV_construction} derives the chain-level $A_\infty$ structure from BV quantization.
Section~\ref{sec:FM_calculus} develops FM boundary calculus, proving the $A_\infty$ relations via Stokes' theorem.
Section~\ref{sec:PVA_descent} establishes cohomological descent to PVAs.
Section~\ref{sec:chiral_hochschild_expanded} relates bulk observables to chiral Hochschild cochains.
Section~\ref{sec:bar_cobar} develops the filtered chiral bar--cobar adjunction.
Section~\ref{sec:examples} presents three worked examples.
Appendix~\ref{app:FM_Stokes} proves orientation conventions, residue formulas, and Arnold relations.

\subsection{Relation to Prior Work}

This manuscript proves foundational results for structures appearing in the physics literature:

\begin{itemize}
\item \textbf{Costello--Gwilliam \cite{CG17}}: Factorization algebras and BV quantization supply the framework; we extract the $A_\infty$ chiral structure implicit in their constructions.

\item \textbf{Costello--Dimofte--Gaiotto \cite{CDG20}}: Boundary chiral algebras and bulk PVAs are constructed physically; proofs and the chain-level $A_\infty$ structure appear here.

\item \textbf{Khan--Zeng \cite{KZ25}}: PVA axioms and the 3d HT Poisson sigma model supply explicit formulas; we derive the $\lambda$-Jacobi identity from $A_\infty$ relations.

\item \textbf{Gaiotto--Kulp--Wu \cite{GKW25}}: Higher operations and configuration space techniques; we complete the proofs via FM calculus.

\item \textbf{Gui--Li--Zeng \cite{GLZ22}}: Quadratic chiral duality is our benchmark; we extend to the non-quadratic $A_\infty$ setting.

\item \textbf{Tamarkin \cite{Tam02}}: Chiral Hochschild theory and $*$-Lie deformations provide the algebraic framework; we give a geometric realization in the HT context.
\end{itemize}

\subsection{Acknowledgments}

This work is informed by conversations with Kevin Terng and draws on talks by Davide Gaiotto.
